\section{Del diseño de flujos al diseño de contexto}

La metodología de producto más importante de este episodio es el paso del flow design al context design. Durante los últimos 20 años, el paradigma básico de los productos de internet ha sido este: el gerente de producto define de antemano la ruta del usuario, las entradas y las salidas; el diseñador dibuja la interfaz, y el ingeniero implementa el flujo. En los productos AI Native, la entrada del usuario es abierta y la salida del modelo también lo es. Sobre todo en la era de los Agent, el modelo además llama herramientas, genera archivos y ejecuta acciones, por lo que el flujo no puede enumerarse de forma exhaustiva.

\begin{figure}[H]
\centering
\includegraphics[width=0.96\textwidth]{flow-to-context-design.png}
\caption{Del diseño de flujos al diseño de contexto: de grados de libertad limitados a una entrada abierta y una salida abierta.}
\end{figure}

\begin{knowledgebox}{Lectura de la figura: en AI Native cambió el objeto del diseño}
A la izquierda, el diseño de flujos se ocupa de botones, páginas, embudos y salidas fijas; a la derecha, el diseño de contexto se ocupa del prompt, el context, las herramientas, la combinación de modelos, las restricciones y la retroalimentación. El flujo del producto sigue siendo importante, pero debe deducirse a partir de las capacidades del modelo y del efecto del contexto, en lugar de dibujar primero el flujo completo y esperar que el modelo se ajuste a él.
\end{knowledgebox}

\subsection{Por qué falla la colaboración tradicional}

La figura anterior muestra la diferencia de paradigma; aquí se aplica a la colaboración del equipo. AI Native no es un cambio de mentalidad exclusivo del gerente de producto, sino un cambio en el orden en que colaboran los equipos de producto, diseño, ingeniería y modelos.

La colaboración tradicional de un equipo es lineal: el gerente de producto escribe los requisitos, el diseñador hace los diseños y el ingeniero los construye. En un producto AI Native, la parte intermedia más importante puede ser el context, que el usuario no ve: qué materiales se le dan al modelo, cómo se organiza la tarea, qué herramientas se llaman, cómo se evalúa la salida y cómo se incorpora la retroalimentación del usuario a la siguiente ronda. Si estos elementos no se hacen funcionar primero de principio a fin, cuanto más refinado sea el flujo de la UI, más equivocado puede estar.

\begin{warningbox}{El riesgo de empezar por el flujo}
Si primero se diseñan por completo las páginas, los botones y las rutas del usuario, y después se conecta el modelo, es muy probable descubrir que la salida real del modelo no coincide con el flujo previsto. Un producto AI Native debe validar primero el ciclo cerrado de capacidades y después diseñar el flujo externo.
\end{warningbox}

\subsection{Resumen de la sección}

El paso del diseño de flujos al diseño de contexto es la brecha conceptual decisiva para el gerente de producto de internet que se mueve hacia los productos de IA. El gerente de producto ya no es solo el diseñador de la ruta del usuario, sino que debe convertirse en quien organiza las capacidades del modelo, el contexto, la retroalimentación y el modelo mental del usuario.
